% ===================== TITLE PAGE =====================
\begin{titlepage}
\thispagestyle{empty}
\vspace*{28mm}
{\color{rule}\hrule height 1.2pt}
\vspace{10mm}

{\sffamily\bfseries\fontsize{27}{31}\selectfont\color{slate}
Shared-Ride Platform\par}
\vspace{5mm}
{\sffamily\fontsize{15}{19}\selectfont\color{ink}
Scheduled shared transport for Alexandria, Egypt\par}

\vspace{9mm}
{\color{rule}\hrule height 0.4pt}
\vspace{9mm}

{\large\color{muted}
Information document describing the product design, operating model,\\
revenue mechanics, evidence base and known risks.\par}

\vfill

\begin{tabular}{@{}L{42mm}L{80mm}@{}}
\thead{Document} & Product and operating design \\
\thead{Version} & 1.0 \\
\thead{Date} & August 2026 \\
\thead{Status} & Pre-development. Specification complete. \\
\thead{Prepared from} & Master Specification v3.0 (167 recorded decisions) \\
\end{tabular}

\vspace{10mm}
{\color{rule}\hrule height 0.4pt}
\vspace{4mm}
{\footnotesize\color{muted}
This document describes work completed to date and design intent. It contains no revenue history,
no user numbers and no financial projections, because none exist. Figures drawn from third-party
research are attributed at the point of use and listed in Appendix~B.\par}
\end{titlepage}

% ===================== CONTENTS =====================
\pagenumbering{roman}
\setcounter{page}{2}
\microtypesetup{protrusion=false}
\tableofcontents
\microtypesetup{protrusion=true}
\cleardoublepage
\pagenumbering{arabic}
\setcounter{page}{1}

% ===================== 1. SUMMARY =====================
\chapter{Summary}

\lead{A scheduled shared-transport service in which an operator defines fixed routes and departure
times, drivers claim the departures they wish to work, and passengers buy a seat at one flat price
per route.}

\section{What the service is}

A passenger opens the application, selects a route, selects a boarding point on that route, and
selects a departure time. They pay a single fixed price regardless of how far along the route they
travel, board by presenting a code, and leave the vehicle wherever they choose along the way.

The operator defines the route network: which stops exist, in what order, at what price, and during
which hours. The operator also publishes a grid of departure times. Drivers do not design routes or
set prices; they select the departures they want to run, in two steps.

In structure this is a microbus network operated through software. The organising idea is not to
replace the informal shared transport that Egyptians already use, but to make it schedulable,
priced in advance, and accountable.

\section{The four design choices that shape everything else}

\begin{keybox}
\panelhead{One} \textbf{Passengers buy a route, not a destination.}\\
A ticket grants travel on a route. The passenger boards at a defined stop and alights wherever they
wish. No destination is declared at purchase. This mirrors existing microbus behaviour and removes
a large amount of complexity from pricing and matching.

\vspace{5pt}
\panelhead{Two} \textbf{One flat fare per route.}\\
The same price applies whether the passenger travels two stops or twelve. The price is fixed when
the ticket is bought and does not change afterwards. There is no demand-based pricing.

\vspace{5pt}
\panelhead{Three} \textbf{Passengers walk to a stop.}\\
Boarding happens at defined meeting points rather than at the passenger's door. Published industry
figures indicate that walk-to-a-point services sustain discounts of roughly 25--40\,\% against a
private ride, while door-to-door shared services sustain roughly 15--22\,\%. The walk is what makes
the price difference possible.

\vspace{5pt}
\panelhead{Four} \textbf{The operator schedules; drivers choose when to work.}\\
Routes, fares and timetables are centrally defined. Drivers claim published departure slots. They
carry no planning burden and take no pricing decisions.
\end{keybox}

\section{Current status}

\begin{center}
\begin{tabularx}{\textwidth}{@{}L{52mm}Y@{}}
\toprule
\thead{Design specification} & Complete. 167 recorded decisions, 22 chapters, 75 defined screens, 65 configuration parameters, 34 system invariants. \\
\thead{Software} & Not started. No code has been written. \\
\thead{Operations} & Not started. No routes mapped, no drivers recruited. \\
\thead{Revenue} & None. The service does not operate. \\
\thead{Regulatory position} & Referred to legal counsel. Not assessed in this document. \\
\bottomrule
\end{tabularx}
\end{center}

\vspace{4pt}
This document therefore describes a design, not a business in operation. Where evidence exists it is
external research, cited at the point of use. Where something is unknown, this document says so.

\section{Why the design took this shape}

The specification began as an on-demand pooled-ride concept and changed direction twice, in both
cases because the available evidence pointed elsewhere.

\begin{enumerate}[label=\textbf{\arabic*.}]
\item \textbf{Towards scheduled routes.} A 2026 evaluation framework covering fourteen studies
      published between 2007 and 2025 reports that fixed-route transit outperforms
      demand-responsive transit on cost per passenger, schedule adherence and ridership wherever
      demand is dense and regular, while demand-responsive services perform better in low-density
      areas. Published cost comparisons are wide: one North American agency recorded roughly
      \$72 per passenger on a flexible service against roughly \$25 on the fixed route it replaced.
      The intended launch context --- university corridors at peak commuting hours --- is the
      condition under which fixed routes perform best.
\item \textbf{Towards recurring revenue.} Swvl, founded in Cairo in 2017 with a comparable
      consumer-facing model, moved away from consumer mass-market operation and reported reaching
      profitability in 2025 on business and government contracts with a high proportion of
      recurring revenue. Subscriptions are consequently central to this design rather than an
      accessory to it.
\end{enumerate}

\srcnote{full citations in Appendix~B.}
